\chapter{Requerimientos del Subsistema SAI}
\label{ch:requerimientos}

Este capítulo define los requerimientos formales del \textbf{Subsistema de
Super-Resolución basada en Inteligencia Artificial (SAI)} del proyecto INTISAT.
Los requerimientos se formulan siguiendo principios de ingeniería de sistemas
ECSS/NASA y se enfocan en funcionalidad, desempeño computacional, interfaces,
operación y limitaciones del subsistema.

Es importante destacar que los requerimientos del SAI se formulan como
\textbf{requerimientos de postprocesamiento}, y no como requerimientos de
adquisición o generación de información científica primaria.

% -------------------------------------------------------------------
\section{Clasificación de métodos de verificación}

Los métodos de verificación utilizados se definen como:

\begin{itemize}
    \item \textbf{T} -- Test: ejecución controlada del software y evaluación de
    resultados.
    \item \textbf{A} -- Analysis: análisis cuantitativo o cualitativo de salidas.
    \item \textbf{I} -- Inspection: inspección visual de resultados procesados.
    \item \textbf{R} -- Review of Design: revisión de arquitectura, código o
    documentación.
\end{itemize}

% -------------------------------------------------------------------
\section{Requerimientos funcionales (FR)}

\begin{longtable}{p{0.12\textwidth} p{0.68\textwidth} p{0.12\textwidth}}
\toprule
\textbf{ID} & \textbf{Descripción} & \textbf{Verificación} \\
\midrule
\endfirsthead
\toprule
\textbf{ID} & \textbf{Descripción} & \textbf{Verificación} \\
\midrule
\endhead
FR-SAI-01 & El subsistema SAI debe recibir como entrada imágenes digitales
previamente adquiridas por el CMOS Microscopy Payload. & R \\
FR-SAI-02 & El subsistema debe aplicar modelos de super-resolución basados en
aprendizaje profundo como etapa de postprocesamiento. & T \\
FR-SAI-03 & El sistema debe generar una imagen procesada sin modificar ni
sobrescribir la imagen original de entrada. & T, I \\
FR-SAI-04 & El subsistema debe permitir seleccionar el modelo de
super-resolución a emplear. & T \\
FR-SAI-05 & El sistema debe conservar la trazabilidad entre imagen de entrada y
salida procesada. & R \\
\bottomrule
\end{longtable}

% -------------------------------------------------------------------
\section[Requerimientos de desempeño (PR)]{Requerimientos de desempeño computacional (PR)}

\begin{longtable}{p{0.12\textwidth} p{0.68\textwidth} p{0.12\textwidth}}
\toprule
\textbf{ID} & \textbf{Descripción} & \textbf{Verificación} \\
\midrule
\endfirsthead
\toprule
\textbf{ID} & \textbf{Descripción} & \textbf{Verificación} \\
\midrule
\endhead
PR-SAI-01 & El subsistema debe ejecutar inferencia de super-resolución en tiempo
finito sobre imágenes individuales en entorno de laboratorio o en tierra. & T \\
PR-SAI-02 & El tiempo de procesamiento no debe impedir el análisis posterior de
los datos adquiridos. & A \\
PR-SAI-03 & El consumo de memoria del subsistema debe ser compatible con la
plataforma de ejecución seleccionada. & A, R \\
PR-SAI-04 & El subsistema no debe introducir artefactos visuales evidentes que
distorsionen las estructuras presentes en la imagen original. & I, A \\
\bottomrule
\end{longtable}

% -------------------------------------------------------------------
\section{Requerimientos de interfaz (IR)}

\begin{longtable}{p{0.12\textwidth} p{0.68\textwidth} p{0.12\textwidth}}
\toprule
\textbf{ID} & \textbf{Descripción} & \textbf{Verificación} \\
\midrule
\endfirsthead
\toprule
\textbf{ID} & \textbf{Descripción} & \textbf{Verificación} \\
\midrule
\endhead
IR-SAI-01 & El subsistema debe aceptar imágenes en formatos estándar sin
compresión (PNG o TIFF). & T \\
IR-SAI-02 & El subsistema debe generar salidas en un formato compatible con
herramientas de análisis en tierra. & T \\
IR-SAI-03 & El subsistema debe mantener compatibilidad con los metadatos
generados por el CMOSM. & R \\
\bottomrule
\end{longtable}

% -------------------------------------------------------------------
\section{Requerimientos operacionales (OR)}

\begin{longtable}{p{0.12\textwidth} p{0.68\textwidth} p{0.12\textwidth}}
\toprule
\textbf{ID} & \textbf{Descripción} & \textbf{Verificación} \\
\midrule
\endfirsthead
\toprule
\textbf{ID} & \textbf{Descripción} & \textbf{Verificación} \\
\midrule
\endhead
OR-SAI-01 & El subsistema SAI debe operar de manera independiente a la captura de
datos del CMOSM. & R \\
OR-SAI-02 & El subsistema debe permitir la ejecución manual del procesamiento por
parte del operador. & T \\
OR-SAI-03 & El uso del subsistema no debe ser obligatorio para la interpretación
de las imágenes originales. & R \\
\bottomrule
\end{longtable}

% -------------------------------------------------------------------
\section{Limitaciones explícitas del subsistema}

Con el fin de evitar interpretaciones erróneas, se establecen las siguientes
limitaciones explícitas del subsistema SAI:

\begin{itemize}
    \item El SAI no genera información física nueva ni incrementa la resolución
    óptica real del sistema.
    \item Los resultados del procesamiento no deben utilizarse como única base
    para conclusiones científicas cuantitativas.
    \item El subsistema no sustituye procesos de calibración física del sistema
    de adquisición.
\end{itemize}

% -------------------------------------------------------------------
\section{Resumen}

Los requerimientos definidos en este capítulo establecen un marco claro y
conservador para el desarrollo e integración del subsistema SAI. La formulación
adoptada garantiza coherencia con el CMOS Microscopy Payload y limita
explícitamente el alcance del uso de técnicas de super-resolución dentro del
proyecto INTISAT.
